% !TEX root = ../main.tex

\chapter{相关技术与理论基础}

本章介绍本文研究所涉及的相关技术与理论基础。首先，对代码生成任务的定义、输入输出形式和技术难点进行说明；其次，介绍预训练代码大模型和监督微调方法，为本文采用 Qwen2.5-Coder-1.5B 进行 SFT 训练提供背景；然后，对强化学习和 PPO 算法的核心思想进行阐述；最后，介绍基于执行反馈的代码生成优化方法，以及本文实验中使用的数据集和评价指标。

\section{代码生成任务概述}

代码生成是指模型根据给定的自然语言描述、函数签名、上下文代码或接口约束，自动生成满足需求的程序代码。与一般自然语言生成任务不同，代码生成的输出不仅需要在语法层面符合编程语言规范，还需要在语义层面满足题目要求，并能够在给定测试环境中正确运行。对于函数级代码生成任务而言，输入通常包含自然语言问题描述和函数入口信息，输出则是一个完整函数实现。模型生成的代码需要与题目提供的测试用例交互，并返回符合预期的结果。

形式化地，可以将代码生成任务表示为条件生成问题。设输入问题描述为 $x$，目标代码序列为 $y=(y_1,y_2,\ldots,y_T)$，自回归语言模型通过最大化条件概率生成代码：

\begin{equation}
  p_{\theta}(y|x)=\prod_{t=1}^{T}p_{\theta}(y_t|x,y_{<t}),
\end{equation}

其中，$\theta$ 表示模型参数，$y_{<t}$ 表示当前 token 之前已经生成的代码序列。模型在推理阶段根据该条件概率分布逐步生成代码，最终得到完整程序。对于预训练语言模型而言，该建模方式与自然语言生成类似，但代码生成任务具有更强的结构约束。

代码生成任务的难点主要体现在三个方面。第一，代码具有严格的语法约束。模型输出中任何缺失缩进、括号不匹配、变量未定义或函数名错误，都可能导致程序无法执行。第二，代码具有明确的语义约束。某段代码即使语法正确，也可能因为算法逻辑错误、边界条件遗漏或返回类型不符而无法通过测试。第三，代码具有多解性。同一问题往往存在多种正确实现方式，基于文本相似度的评价并不能充分反映程序功能正确性。因此，代码生成评价通常更关注执行结果，而不是仅比较生成代码与参考答案的字面相似度。

本文关注的是函数级 Python 代码生成任务。给定问题描述和入口函数名称，模型需要生成对应函数实现。生成代码会被放入测试环境中执行，并根据测试用例通过情况计算奖励或评价指标。这一任务形式与 HumanEval 和 MBPP 等常用代码生成基准保持一致，也便于将执行结果作为强化学习反馈。

\section{预训练代码大模型}

预训练代码大模型是将大规模语言模型技术应用于程序语言建模的结果。Transformer 架构通过自注意力机制对序列中不同位置的 token 建立关联，使模型能够捕捉长距离依赖关系\cite{Vaswani2017Attention}。在自然语言和代码语料上进行大规模预训练后，模型能够学习词法结构、语法模式、库函数调用、常见算法模板以及自然语言描述与代码之间的对应关系。

从训练目标看，代码大模型通常采用自回归语言建模方式，即给定前文 token 预测下一个 token。对于代码数据而言，这一目标能够促使模型学习代码片段的连续结构。例如，在 Python 函数中，模型需要根据函数签名、变量名和上下文推断后续控制流、表达式和返回值。随着模型规模和训练数据规模扩大，预训练代码模型逐渐具备从自然语言问题描述直接生成代码的能力。

Codex 的研究系统展示了大语言模型在代码生成任务中的能力，并提出 HumanEval 作为重要评价基准\cite{Chen2021Codex}。后续 Code Llama、DeepSeek-Coder 等开源代码模型进一步推动了代码大模型的发展\cite{Roziere2023CodeLlama,Guo2024DeepSeekCoder}。这些模型通常在大规模代码仓库、技术文档和自然语言语料上训练，能够处理代码补全、代码解释、代码翻译和自然语言到代码生成等任务。

尽管预训练代码大模型已经具备较强的通用能力，但在具体任务场景中仍然需要进一步适配。一方面，不同数据集的 prompt 格式、函数签名和输出要求存在差异；另一方面，预训练模型生成代码时可能输出解释性文本、Markdown 代码块或额外测试代码，导致评测脚本难以直接执行。因此，在本文实验中，预训练代码模型既作为 Base 进行直接评测，也作为 SFT 和 PPO 的初始化模型。通过比较 Base、SFT 和 PPO，可以观察监督微调与执行反馈强化学习分别对模型行为产生的影响。

本文选用 Qwen2.5-Coder-1.5B 作为基座模型。该模型属于代码专用因果语言模型，本身已经具备较强的代码生成能力。选择 1.5B 规模模型的原因在于，它能够在单张消费级 GPU 上完成 LoRA 微调和 PPO 实验，兼顾实验可行性和代码生成能力。由于本文研究重点是执行反馈训练流程与奖励设计，而不是从零训练大模型，因此采用强代码基座模型作为实验起点更符合实际应用场景。

\section{监督微调方法}

监督微调是指在预训练模型基础上，使用标注数据进一步训练模型，使其适配特定任务格式和输出分布。在代码生成任务中，监督微调数据通常由输入 prompt 和目标代码 response 构成。模型训练目标是最大化参考答案在输入条件下的似然，即最小化 token 级交叉熵损失：

\begin{equation}
  \mathcal{L}_{\mathrm{SFT}}(\theta)
  = -\sum_{t=1}^{T}\log p_{\theta}(y_t|x,y_{<t}).
\end{equation}

其中，$x$ 表示输入问题描述，$y_t$ 表示参考代码序列中的第 $t$ 个 token。通过该目标，模型学习在给定问题描述后生成参考实现。监督微调的优点是训练过程稳定、实现简单，并且能够快速使模型适应指定输出格式。因此，在许多语言模型后训练流程中，SFT 通常作为强化学习之前的热启动阶段。

在代码生成场景中，SFT 能够带来两个直接作用。第一，它可以使模型更稳定地输出函数实现，而不是混入过多解释性文本。第二，它可以使模型学习目标数据集中的函数命名、解题风格和常见任务形式。对于本文使用的 MBPP 数据而言，SFT 训练样本由自然语言问题描述和标准答案构成，模型通过训练学习从问题描述到函数代码的映射。

但 SFT 也存在局限。由于训练目标是模仿参考答案，模型优化的是 token 级相似性，而不是程序执行正确性。代码问题通常存在多种正确实现方式，某种实现与参考答案 token 差异较大，并不意味着它是错误的；相反，与参考答案局部相似的实现也可能因为边界条件错误而失败。因此，SFT 的优化目标与代码生成最终评价目标之间存在不一致。此外，当 SFT 数据规模较小或分布较窄时，模型可能过度适配训练数据形式，导致部分通用能力下降。本文实验中也将 SFT 作为独立基线，用于分析小样本 SFT 对强代码基座模型的影响。

为了降低训练成本，本文采用 LoRA 进行参数高效微调。LoRA 的核心思想是在冻结原模型大部分参数的情况下，为部分线性变换层引入低秩可训练矩阵，使模型能够通过少量新增参数适配下游任务。设原始权重矩阵为 $W$，LoRA 将参数更新表示为低秩分解：

\begin{equation}
  W' = W + \Delta W = W + BA,
\end{equation}

其中，$A$ 和 $B$ 为可训练低秩矩阵。该方法能够显著降低显存需求和训练参数量，适合在单张 GPU 上完成代码模型微调。本文在 Qwen2.5-Coder-1.5B 的多个注意力和前馈模块中应用 LoRA，以获得可用于后续 PPO 优化的 SFT adapter。

\section{强化学习基础}

强化学习研究智能体如何通过与环境交互学习最优策略。一个典型强化学习问题可以用马尔可夫决策过程表示，包括状态空间 $\mathcal{S}$、动作空间 $\mathcal{A}$、状态转移概率、奖励函数 $r$ 和折扣因子 $\gamma$。在每一步，智能体根据当前状态 $s_t$ 选择动作 $a_t$，环境返回奖励 $r_t$ 并转移到下一个状态 $s_{t+1}$。强化学习的目标是学习策略 $\pi_{\theta}(a|s)$，最大化期望累计回报：

\begin{equation}
  J(\theta)=\mathbb{E}_{\pi_{\theta}}\left[\sum_{t=0}^{T}\gamma^t r_t\right].
\end{equation}

在语言模型生成任务中，可以将输入 prompt 视为状态的一部分，将模型逐步生成 token 的过程视为动作序列。对于代码生成任务，模型完成一次代码生成后，环境通过执行测试返回奖励。因此，代码生成可以被建模为一种序列决策问题：模型策略决定生成什么代码，执行环境根据代码质量给出奖励。

价值函数和优势函数是强化学习中的重要概念。状态价值函数 $V^{\pi}(s)$ 表示在状态 $s$ 下按照策略 $\pi$ 行动所能获得的期望回报；动作价值函数 $Q^{\pi}(s,a)$ 表示在状态 $s$ 下采取动作 $a$ 后继续按照策略 $\pi$ 行动的期望回报。优势函数用于衡量某一动作相对于当前状态平均水平的好坏：

\begin{equation}
  A^{\pi}(s,a)=Q^{\pi}(s,a)-V^{\pi}(s).
\end{equation}

在代码生成实验中，完整程序生成完成后才能执行测试，因此奖励通常是 episode 级别的延迟反馈。本文采用单步 episode 近似，将一次完整代码生成视为一个决策样本，基于执行结果计算 reward 和 advantage。虽然这种实现比完整多步轨迹 PPO 更简化，但足以验证执行反馈能否对代码生成策略产生有效优化。

\section{PPO 算法原理}

PPO 是一种常用的策略梯度强化学习算法，其核心目标是在提升策略收益的同时限制策略更新幅度，避免新策略偏离旧策略过远而导致训练不稳定\cite{Schulman2017PPO}。在语言模型后训练中，PPO 被广泛用于根据奖励模型或外部反馈优化生成策略。

PPO 首先定义新旧策略概率比：

\begin{equation}
  r_t(\theta)=\frac{\pi_{\theta}(a_t|s_t)}{\pi_{\theta_{\mathrm{old}}}(a_t|s_t)}.
\end{equation}

如果直接最大化 $r_t(\theta)A_t$，当优势估计较大时，策略可能发生过度更新。为缓解这一问题，PPO 引入 clipped objective：

\begin{equation}
  \mathcal{L}^{\mathrm{CLIP}}(\theta)
  = \mathbb{E}_t
  \left[
  \min
  \left(
  r_t(\theta)A_t,
  \mathrm{clip}(r_t(\theta),1-\epsilon,1+\epsilon)A_t
  \right)
  \right],
\end{equation}

其中，$\epsilon$ 为裁剪范围超参数。该目标限制新策略相对旧策略的变化幅度，使训练过程更加稳定。

在实际训练中，PPO 通常还会引入值函数损失和 KL 约束。值函数用于估计当前状态的期望回报，从而辅助计算优势；KL 约束则用于限制当前策略与参考策略之间的差异，防止模型为了追求短期奖励而破坏原有语言建模能力。一个常见的总损失形式可以写为：

\begin{equation}
  \mathcal{L}
  = \mathcal{L}_{\mathrm{policy}}
  + c_v\mathcal{L}_{\mathrm{value}}
  + c_{kl}\mathcal{L}_{\mathrm{KL}},
\end{equation}

其中，$c_v$ 和 $c_{kl}$ 分别表示 value loss 和 KL loss 的权重。

本文实现的 PPO 流程包含 clipped policy loss、scalar value head、单步优势估计和 KL penalty。由于本文实验规模较小，PPO 被设计为轻量级单步优化框架：模型先基于 SFT adapter 生成代码，执行环境计算 reward，再用这些 rollout 样本进行一次或多次 mini-batch 更新。该设计能够在有限算力条件下验证执行反馈强化学习的有效性。

\section{基于执行反馈的代码生成优化}

基于执行反馈的代码生成优化利用程序语言可执行的特点，将模型生成结果放入真实或受控执行环境中运行，并根据执行结果构造反馈信号。与自然语言任务中依赖人工偏好或文本相似度不同，代码任务可以直接通过解释器、编译器和单元测试判断输出质量。这使得代码生成任务天然适合引入外部环境反馈。

执行反馈通常包括以下几类信息。第一，可执行性反馈，即代码是否能够被解释器或编译器成功解析和运行。第二，接口一致性反馈，即生成代码是否包含预期入口函数，函数参数和返回形式是否符合要求。第三，测试通过反馈，即生成代码通过了多少单元测试。第四，错误类型反馈，例如语法错误、运行时异常、断言失败或超时。这些反馈既可以用于最终评测，也可以进一步转化为强化学习奖励。

在代码生成强化学习中，奖励函数设计非常关键。如果只在全部测试通过时给予奖励，模型会面对较稀疏的反馈，大量接近正确但未完全通过的代码得不到有效区分。如果奖励过于复杂，例如过度依赖代码结构相似度，又可能使模型优化代理指标而非真正的功能正确性。因此，执行反馈奖励需要在密度和目标一致性之间取得平衡。

本文围绕这一问题设计了多种 reward。原始 reward 主要基于测试通过比例；reward\_v2 在此基础上加入代码可执行奖励和全部通过奖励，使反馈更加密集；reward\_v3 和 reward\_v3b 进一步尝试引入 AST 相似度和函数签名匹配奖励，用于分析结构奖励是否能够提升代码生成质量。第四章实验将表明，在本文小规模数据和单步 PPO 设置下，执行正确性相关奖励比结构相似度奖励更稳定有效。

\section{代码生成数据集与评价指标}

本文实验主要使用 MBPP 和 HumanEval 两个代码生成数据集。MBPP 是面向基础 Python 编程问题的数据集，题目通常由自然语言描述、参考答案和若干测试用例构成。该数据集问题规模相对较小，适合用于监督微调和强化学习 rollout。HumanEval 是函数级代码生成评价基准，题目通常包含函数签名、文档字符串和隐藏测试逻辑，重点考察模型根据问题描述生成正确函数实现的能力\cite{Chen2021Codex}。

在本文实验中，MBPP 被划分为训练集和验证集，用于 SFT 训练、PPO rollout 和验证集评估；HumanEval 作为最终测试集，用于观察模型在跨数据集场景下的泛化能力。为了保证数据可用性，本文将原始数据统一转换为包含 task\_id、prompt、canonical\_solution、entry\_point 和 test\_code 等字段的格式，并对标准答案进行 runtime 复核，确保测试代码能够正确执行。

代码生成任务最常用的指标之一是 Pass@1。对于每个问题，模型生成一个候选答案，如果该答案通过全部测试，则记为通过；Pass@1 表示所有问题中一次生成即通过的比例：

\begin{equation}
  \mathrm{Pass@1}=\frac{N_{\mathrm{pass}}}{N_{\mathrm{total}}},
\end{equation}

其中，$N_{\mathrm{pass}}$ 表示通过全部测试的样本数，$N_{\mathrm{total}}$ 表示总样本数。Pass@1 能够直接反映模型在单次生成场景下的代码正确率，也是本文比较 Base、SFT 和 PPO 的核心指标。

除 Pass@1 外，本文还记录 average reward 和 generation execution failures。average reward 表示所有样本执行反馈奖励的平均值，可以反映模型整体输出质量；generation execution failures 表示生成代码在执行阶段发生失败的数量，可以衡量模型输出的可执行稳定性。对于 PPO 训练过程，还记录 policy loss、value loss、KL loss 和平均策略概率比等指标，用于分析策略更新是否稳定。

\section{本章小结}

本章介绍了本文研究所需的相关技术与理论基础。首先，代码生成任务被定义为一种以自然语言或函数签名为输入、以可执行程序为输出的条件生成问题，其评价重点是功能正确性而非文本相似性。其次，预训练代码大模型为代码生成提供了强大的初始能力，而监督微调能够使模型适配特定任务格式，但也存在训练目标与执行正确性不一致的问题。随后，本章介绍了强化学习和 PPO 算法的基本原理，说明 PPO 通过裁剪目标和 KL 约束限制策略更新幅度。最后，本章讨论了执行反馈在代码生成优化中的作用，并介绍了本文使用的 MBPP、HumanEval、Pass@1 和 reward 等数据集与评价指标。以上内容为第三章提出基于执行反馈的 PPO 代码生成方法奠定了理论基础。
